\section{Project Background}
\label{sec:1-1}

The project focuses on customer retention in an e-commerce business. Revenue is generated through customer purchases, while customer acquisition and retention both require business resources. The project investigates how Data Science can support the retention of existing customers by predicting churn.

\section{Business Context}
\label{sec:1-2}

The business generates sales revenue from new customers making first purchases and from existing customers making repeat purchases. Existing customers may stop purchasing, which creates a customer churn problem.

\section{Problem Statement}
\label{sec:1-3}

The business needs to increase sales revenue while operating under limited retention and marketing resources. The project therefore focuses on identifying existing customers at risk of churn so that retention resources can be prioritized.

\section{Research Questions}
\label{sec:1-4}

\begin{enumerate}
    \item How can transaction history be transformed into customer-level observations for churn prediction?
    \item How effectively can purchasing behavior predict future churn?
    \item Can churn-risk ranking support prioritization under limited retention capacity, compared with simple rule-based targeting?
    \item Which behavioral characteristics are \emph{associated} with churn? (Associational, not causal.)
\end{enumerate}

\section{Project Objectives}
\label{sec:1-5}

Design and evaluate a customer churn prediction system that supports the retention of existing customers by producing customer-level churn risk and supporting customer prioritization.

\section{Scope and Limitations}
\label{sec:1-6}

\textbf{In scope:} existing-customer retention, churn prediction, customer snapshots, weekly batch scoring, risk ranking and Top-K targeting.

\textbf{Out of scope / not claimed:}
\begin{itemize}
    \item Causal uplift and incremental campaign revenue -- the transaction data contain no randomized treatment/control information.
    \item A live deployment. \textbf{Chapters~\ref{ch:deployment} and~\ref{ch:feedback} describe a proposed operational design}; they are not evaluated on a running system.
    \item Profit analysis based on real margins or costs -- these are not in the data and enter only as stated assumptions.
\end{itemize}

\section{Report Structure}
\label{sec:1-7}

The report follows the 12 chapters defined in this strategy.
